%%=============================================================================
%% Methodologie
%%=============================================================================

\chapter{\IfLanguageName{dutch}{Methodologie}{Methodology}}%
\label{ch:methodologie}

In dit hoofdstuk wordt toegelicht hoe het onderzoek praktisch werd uitgevoerd en hoe de verschillende experimenten werden opgebouwd. De methodologie vertrekt vanuit de centrale onderzoeksvraag: hoe graph-based machine learning kan worden toegepast om security-events automatisch te correleren en te groeperen tot mogelijke aanvalsketens.
\newline
\newline
De uitwerking bestaat uit een opeenvolgend proces waarbij ruwe security-events eerst worden voorbereid en omgezet naar een graafstructuur. Vervolgens wordt een relationeel graphmodel getraind waarna de output gebruikt wordt voor verdere analyse, ranking en threat chaining.
\newline
\newline
De volledige pipeline bestaat uit datasetselectie, preprocessing, grafconstructie, modeltraining, evaluatie en chaining.

\section{Onderzoeksopzet}

Het onderzoek werd opgebouwd rond twee datasets met een verschillende doelstelling.
\newline
\newline
CIC-IDS2018 werd gebruikt als gecontroleerde benchmark voor supervised multiclass classificatie. Deze dataset bevat labels voor verschillende aanvalstypes en laat toe om de performantie van het model te evalueren in een klassieke detectiesetting.
\newline
\newline
LANL werd gebruikt als meer realistische enterprise dataset. Deze bevat flow-, authenticatie-, DNS- en procesdata, maar beschikt niet over volledige labels per event. Hierdoor werd LANL behandeld als een weakly supervised probleem waarbij de nadruk verschuift van pure classificatie naar ranking en chaining.
\newline
\newline
Hoewel dezelfde modelarchitectuur gebruikt werd voor beide experimenten, verschilt de evaluatiecontext:
\begin{itemize}
    \item CIC-IDS2018: supervised multiclass classificatie;
    \item LANL: suspicious event scoring, ranking en threat chaining.
\end{itemize}

\section{Datasetselectie}

\subsection{CIC-IDS2018}

CIC-IDS2018 werd gebruikt als primaire dataset voor multiclass attack detection.
\newline
\newline
De dataset bevat netwerkflows met labels voor verschillende aanvalstypes zoals brute force, DDoS, DoS, infiltratie en SQL-injectie. Hierdoor kan onderzocht worden in welke mate het model verschillende aanvalsklassen van elkaar kan onderscheiden.
\newline
\newline
De data bestaat voornamelijk uit flowgebaseerde netwerkfeatures. Hoewel deze dataset geen volledige bedrijfscontext bevat, laat ze wel toe om relationele patronen tussen netwerkgebeurtenissen te modelleren.

\begin{table}[H]
    \centering
    \caption{Overzicht van de gebruikte CIC-IDS2018 subset}
    \label{tab:cic_dataset}
    
    \begin{tabular}{p{4cm}p{7cm}}
        \toprule
        Kenmerk & Waarde \\
        \midrule
        
        Doel & Multiclass intrusion detection \\
        
        Datatype & Netwerkflows \\
        
        Voornaamste features &
        Bron-IP, doel-IP, poorten, protocol, flowkenmerken, tijdsvensters \\
        
        Aantal klassen &
        Benign, DoS, DDoS, Brute Force, SQL Injection, Infiltration \\
        
        Leercontext &
        Supervised multiclass classificatie \\
        
        Evaluatiemetrics &
        Accuracy, Precision, Recall, Macro F1, Weighted F1 \\
        
        \bottomrule
    \end{tabular}
    
\end{table}

\subsection{LANL}

LANL werd gebruikt voor het onderzoeken van threat chaining.
\newline
\newline
Deze dataset bevat meerdere types enterprise-events zoals authenticatie-, flow-, DNS- en procesdata. Hierdoor is de dataset geschikter voor het modelleren van temporele en relationele verbanden tussen events.
\newline
\newline
Omdat geen volledige labels beschikbaar zijn, werd gebruik gemaakt van weak supervision op basis van redteam-data. Events die overlap vertonen met redteam-entiteiten en tijdsvensters worden beschouwd als verdachte kandidaat-events.

\begin{table}[H]
    \centering
    \caption{Overzicht van de gebruikte LANL eventdata}
    \label{tab:lanl_dataset}
    
    \begin{tabular}{p{4cm}p{7cm}}
        \toprule
        Kenmerk & Waarde \\
        \midrule
        
        Doel &
        Threat chaining en suspicious event ranking \\
        
        Eventtypes &
        Authenticatie, flow-events, DNS-events en procesdata \\
        
        Labeltype &
        Weak supervision op basis van redteam-overlap \\
        
        Relaties &
        Broncomputer, doelcomputer, tijd, eventtype, poorten \\
        
        Uitdagingen &
        Class imbalance en beperkte labels \\
        
        Evaluatiemetrics &
        Precision@K, Recall@K, Hit@K, chain count \\
        
        \bottomrule
    \end{tabular}
    
\end{table}

\begin{listing}[H]
    \caption{Weak-label toekenning binnen \texttt{map\_lanl\_to\_generic\_schema}.}
    \begin{minted}{python}
        labels = np.zeros(len(events), dtype=np.int64)
        w = np.timedelta64(int(window_sec), 's')
        ts = events['timestamp'].values.astype('datetime64[ns]')
        se = events['src_entity'].astype(str).values
        de = events['dst_entity'].astype(str).values
        red_times = redteam['timestamp'].values.astype('datetime64[ns]')
        red_src = redteam['src_computer'].values
        red_dst = redteam['dst_computer'].values
        
        order = np.argsort(ts)
        ts_s = ts[order]
        mode = str(weak_label_mode).lower().strip()
        
        for i in range(len(redteam)):
        t0, t1 = red_times[i] - w, red_times[i] + w
        lo = int(np.searchsorted(ts_s, t0, side='left'))
        hi = int(np.searchsorted(ts_s, t1, side='right'))
        if lo >= hi:
        continue
        idx = order[lo:hi]
        if mode == 'strict':
        m = (
        (se[idx] == red_src[i]) & (de[idx] == red_dst[i])
        ) | (
        (se[idx] == red_dst[i]) & (de[idx] == red_src[i])
        )
        else:
        m = (
        (se[idx] == red_src[i]) | (de[idx] == red_src[i])
        | (se[idx] == red_dst[i]) | (de[idx] == red_dst[i])
        )
        labels[idx[m]] = 1
    \end{minted}
\end{listing}
    

\section{Graphconstructie}

Na de selectie van de datasets werden de ruwe events omgezet naar een graafstructuur.
\newline
\newline
Events en entiteiten worden voorgesteld als nodes, terwijl relaties tussen deze entiteiten worden gemodelleerd als edges.
\newline
\newline
Voor CIC-IDS2018 werden netwerkflows gemodelleerd als eventnodes. Daarnaast werden contextuele entiteiten zoals IP-adressen, poorten, protocollen en tijdsvensters gebruikt om relationele verbanden op te bouwen.
\newline
\newline
Voor LANL werden flow-, authenticatie-, DNS- en proces-events eerst omgezet naar een generiek eventschema. Vervolgens werden relaties gelegd op basis van:
\begin{itemize}
    \item bron- en doelcomputer;
    \item poorten;
    \item eventtype;
    \item temporele nabijheid.
\end{itemize}

Hierdoor ontstaat één gezamenlijke heterogene graaf waarin relationele context expliciet aanwezig blijft.

\section{Preprocessing en leakage-preventie}

Bij graph-based machine learning bestaat het risico dat informatie onbedoeld van train naar test vloeit via gedeelde relaties.
\newline
\newline
Om dit te vermijden werd gewerkt met gescheiden train-, validation- en testdelen.
\newline
\newline
Waar mogelijk werd een tijdsgebaseerde split toegepast zodat het model leert op oudere events en geëvalueerd wordt op latere gebeurtenissen.

\begin{listing}[H]
    \begin{minted}{python}
        def chronological_split(
        df,
        time_col,
        train_ratio,
        val_ratio
        ):
        
        df = df.sort_values(
        time_col
        )
        
        n = len(df)
        
        train_end = int(
        n * train_ratio
        )
        
        val_end = int(
        n *
        (
        train_ratio +
        val_ratio
        )
        )
        
        train_df = df.iloc[
        :train_end
        ]
        
        val_df = df.iloc[
        train_end:val_end
        ]
        
        test_df = df.iloc[
        val_end:
        ]
        
        return (
        train_df,
        val_df,
        test_df
        )
    \end{minted}
    
    \caption[Tijdsgebaseerde datasplit]
    {Opsplitsing van events in train-, validation- en testdata.}
    
    \label{code:chronological-split}
\end{listing}

Daarnaast werd bij CIC-IDS2018 gebruik gemaakt van class-aware sampling om dominante klassen minder zwaar te laten doorwegen.
\newline
\newline
Bij LANL werd balancing enkel toegepast op trainingsdata terwijl validation en test hun natuurlijke verdeling behielden.

\section{Modelontwikkeling}

Voor beide experimenten werd gekozen voor een Relational Graph Convolutional Network (R-GCN).
\newline
\newline
Deze keuze werd gemaakt omdat de grafen verschillende node- en relatietypes bevatten.
\newline
\newline
Het model bestaat uit twee relationele graph convolution lagen gevolgd door dropout en een lineaire classificatielaag.

\begin{listing}[H]
    \begin{minted}{python}
        class EventRGCN(nn.Module):
        
        def __init__(
        self,
        in_dim,
        hidden_dim,
        num_relations,
        num_classes
        ):
        
        super().__init__()
        
        self.conv1 = RGCNConv(
        in_dim,
        hidden_dim,
        num_relations
        )
        
        self.conv2 = RGCNConv(
        hidden_dim,
        hidden_dim,
        num_relations
        )
        
        self.classifier =
        nn.Linear(
        hidden_dim,
        num_classes
        )
    \end{minted}
    
    \caption[R-GCN architectuur]
    {Gebruikte R-GCN modelarchitectuur.}
    
    \label{code:rgcn-model}
\end{listing}

Voor CIC-IDS2018 werd het model gebruikt voor multiclass classificatie.
\newline
\newline
Voor LANL werd dezelfde architectuur gebruikt, maar werd de output geïnterpreteerd als verdachtheidsscore voor ranking en chaining.

\section{Threat chaining}

Na de modeltraining werd een bijkomende chainingstap uitgevoerd.
\newline
\newline
Deze chaining gebeurt niet binnen het model zelf, maar als post-processing op de modeloutput.
\newline
\newline
Verdachte events worden eerst geselecteerd op basis van hun modelscore of rankingpositie.
\newline
\newline
Daarna worden events gegroepeerd wanneer zij:
\begin{itemize}
    \item gedeelde entiteiten bevatten;
    \item binnen hetzelfde tijdsvenster vallen;
    \item dezelfde contextuele kenmerken delen.
\end{itemize}

\begin{listing}[H]
    \begin{minted}{python}
        def build_chains(
        events,
        threshold,
        time_window_sec
        ):
        
        suspicious =
        events[
        events[
        "attack_prob"
        ] >= threshold
        ]
        
        suspicious =
        suspicious.sort_values(
        "timestamp"
        )
        
        chains = []
        
        return chains
    \end{minted}
    
    \caption[Threat chaining]
    {Vereenvoudigde chaininglogica.}
    
    \label{code:threat-chaining}
\end{listing}

Deze stap laat toe om losse voorspellingen te groeperen tot samenhangende onderzoekscases.

\section{Evaluatie}

Voor CIC-IDS2018 werden klassieke classificatiemetrics gebruikt:
\begin{itemize}
    \item accuracy;
    \item precision;
    \item recall;
    \item macro F1;
    \item weighted F1.
\end{itemize}

Voor LANL werd bijkomend gebruik gemaakt van rankinggerichte metrics:
\begin{itemize}
    \item Precision@K;
    \item Recall@K;
    \item Hit@K;
    \item chain count;
    \item gemiddelde chainlengte.
\end{itemize}

Deze keuze werd gemaakt omdat weak labels en sterke class imbalance klassieke classificatiemetrics minder representatief maken.
\newline
\newline
Naast classificatie werd daarom ook gekeken naar chainingkwaliteit en bruikbaarheid binnen een operationele context.

\section{Toepassing in een operationele context}

Naast de experimentele datasets werd onderzocht hoe de methode kan worden toegepast op realistische enterprise data.
\newline
\newline
De beschikbare data bestaat uit JSON-structuren afkomstig uit asset management en attack surface management omgevingen.
\newline
\newline
Hierbij kunnen assets, services, softwarecomponenten, gebruikers en dreigingen rechtstreeks worden voorgesteld als nodes binnen een heterogene graaf.
\newline
\newline
Door deze contextuele relaties mee te nemen ontstaat een uitbreidbare structuur die toepasbaar is binnen operationele securityomgevingen.

\chapter{\IfLanguageName{dutch}{Resultaten en Analyse}{Results and Analysis}}%
\label{ch:resultaten}

In dit hoofdstuk worden de resultaten van de uitgevoerde experimenten besproken en geïnterpreteerd in functie van de onderzoeksvraag.
\newline
\newline
De experimenten werden uitgevoerd op twee datasets met een verschillende doelstelling. CIC-IDS2018 werd gebruikt als gecontroleerde benchmark voor supervised multiclass classificatie, terwijl LANL werd gebruikt als weakly supervised enterprise dataset voor ranking en threat chaining.
\newline
\newline
Omdat beide datasets een andere rol vervullen binnen het onderzoek, worden de resultaten niet rechtstreeks met elkaar vergeleken. In plaats daarvan wordt onderzocht welke inzichten elke dataset oplevert voor graph-based security-event analyse.
\newline
\newline
Naast klassieke classificatiemetrics wordt ook gekeken naar chaining en contextuele interpretatie van events. Hierdoor wordt niet enkel de detectieperformantie geëvalueerd, maar ook de praktische bruikbaarheid van de voorgestelde aanpak.

\section{Resultaten op een gelabelde intrusion detection dataset}

De eerste experimenten werden uitgevoerd op CIC-IDS2018. Deze dataset werd gebruikt als benchmark voor supervised multiclass attack detection.
\newline
\newline
Het R-GCN model behaalde een accuracy van 0.8738, een weighted F1-score van 0.8744 en een macro F1-score van 0.7970.
\newline
\newline
De relatief beperkte afstand tussen weighted F1 en macro F1 wijst erop dat het model niet uitsluitend leert op dominante klassen, maar ook minder frequente aanvalstypes gedeeltelijk correct kan onderscheiden.
\newline
\newline
Toch blijft er een verschil zichtbaar. De lagere macro F1-score geeft aan dat zeldzamere aanvalstypes moeilijker te classificeren blijven. Dit is typisch voor intrusion detection datasets waar sommige aanvallen aanzienlijk minder voorkomen dan normaal netwerkverkeer.
\newline
\newline
De resultaten tonen aan dat het graph-based model relevante patronen leert uit de relationele structuur van de data. Door events niet geïsoleerd te behandelen maar expliciet relaties tussen entiteiten op te nemen, kan bijkomende context meegenomen worden tijdens de classificatie.
\newline
\newline
Deze resultaten beantwoorden gedeeltelijk de onderzoeksvraag binnen het oplossingsdomein: graph-based machine learning kan gebruikt worden om security-events relationeel voor te stellen en verschillende aanvalstypes van elkaar te onderscheiden wanneer duidelijke labels beschikbaar zijn.

\section{Resultaten op een enterprise event dataset}

Voor LANL werd dezelfde modelarchitectuur gebruikt, maar in een andere context.
\newline
\newline
In tegenstelling tot CIC-IDS2018 bevat LANL geen volledige labels per event. Daarom werd deze dataset niet behandeld als een klassieke classificatietaak, maar als een weakly supervised suspicious-event detection probleem.
\newline
\newline
Het model behaalde een accuracy van 0.9892 en een weighted F1-score van 0.9839. Op het eerste zicht lijken deze resultaten zeer sterk, maar de macro F1-score bedraagt slechts 0.4973.
\newline
\newline
Dit verschil wordt grotendeels verklaard door de sterke class imbalance binnen de dataset. Het merendeel van de events behoort tot normaal gedrag, terwijl verdachte events slechts een klein deel van de data vormen.
\newline
\newline
Een hoge accuracy betekent in deze context dus niet noodzakelijk dat het model bedreigingen perfect detecteert. Het betekent vooral dat het model de dominante klasse correct herkent.
\newline
\newline
Om die reden werd de evaluatie uitgebreid met ranking- en chaininggerichte metrics. In plaats van elk event afzonderlijk te classificeren, werd onderzocht of verdachte events voldoende hoog gerangschikt worden om bruikbaar te zijn voor verdere analyse.
\newline
\newline
Deze interpretatie sluit beter aan bij een operationele context waarin analisten niet alle events onderzoeken, maar vooral werken met geprioritiseerde lijsten en onderzoekscases.
\newline
\newline
De resultaten tonen dat graph-based machine learning ook bruikbaar blijft wanneer labels beperkt zijn, al verschuift de rol van het model dan van pure classificatie naar ondersteuning van prioritisatie en chaining.

\section{Analyse van chainingresultaten}

Naast individuele eventdetectie werd ook gekeken naar het vermogen van het model om events te groeperen tot samenhangende chains. Dit sluit beter aan bij hoe aanvallen zich in de praktijk manifesteren, namelijk als een reeks onderling gerelateerde stappen.
\newline
\newline
De gegenereerde chains tonen dat het model in staat is om verbanden te leggen tussen events op basis van gedeelde entiteiten en temporele relaties. In plaats van geïsoleerde detecties ontstaat een meer gestructureerd beeld van mogelijke aanvalspaden.
\newline
\newline
De verdeling van de chainlengtes toont dat zowel korte als langere chains worden gevormd. Korte chains kunnen wijzen op geïsoleerde of beperkte activiteit, terwijl langere chains indicaties kunnen zijn van multi-stage aanvallen waarbij meerdere stappen over tijd plaatsvinden.
\newline
\newline
Hoewel de chainingresultaten geen perfecte reconstructie van volledige aanvalspaden garanderen, tonen ze wel aan dat de graph-based aanpak een duidelijke meerwaarde biedt bij het contextualiseren van events. Door relaties expliciet te modelleren, wordt het mogelijk om patronen te identificeren die moeilijk zichtbaar zijn in traditionele, event-gebaseerde analyses.

\section{Interpretatie in een operationele context}

De combinatie van classificatie-, ranking- en chainingresultaten toont aan dat de performantie van een model sterk afhankelijk is van de gekozen evaluatiecontext.
\newline
\newline
Bij datasets met duidelijke labels en een relatief gestructureerde opbouw kan een graph-based model effectief ingezet worden voor multiclass attack detection. In meer realistische datasets, waar labels beperkt of onvolledig zijn, verschuift de focus naar prioritisatie en contextuele analyse.
\newline
\newline
De resultaten tonen dat een graph-based aanpak niet enkel waarde biedt op vlak van detectie, maar ook bij het structureren en interpreteren van security-events. Dit is vooral relevant in omgevingen waar grote hoeveelheden data gegenereerd worden en waar het leggen van verbanden tussen events essentieel is.
\newline
\newline
Door events te modelleren als een graaf en relaties expliciet op te nemen, ontstaat een flexibele representatie die kan aangepast worden aan verschillende types data en use cases. Dit maakt de aanpak toepasbaar binnen uiteenlopende operationele omgevingen waarin security-events geanalyseerd worden.



\section{Samenvattende vergelijking}

De resultaten tonen dat dezelfde graph-based aanpak op verschillende manieren ingezet kan worden, afhankelijk van de beschikbare data en het doel van de analyse.
\newline
\newline
Bij gelabelde datasets ligt de focus op nauwkeurige classificatie van aanvalstypes. Bij meer complexe en minder gestructureerde datasets verschuift de nadruk naar het identificeren van verdachte patronen en het reconstrueren van relaties tussen events.
\newline
\newline
Deze combinatie van detectie, ranking en chaining vormt een belangrijke stap richting een meer contextuele en schaalbare aanpak voor security-event analyse.

\section{Voorbeeld van een gereconstrueerde chain}
\begin{table}[H]
    \centering
    \caption{Voorbeeld van een gereconstrueerde chain}
    \label{tab:example-chain}
    \begin{tabular}{llll}
        \toprule
        Stap & Eventtype & Entiteit & Relatie \\
        \midrule
        1 & Authentication & Host A & Initiële activiteit \\
        2 & DNS & Host A & Zelfde host \\
        3 & Process & Host A & Temporele overlap \\
        4 & Network flow & Host A → Host B & Communicatie \\
        \bottomrule
    \end{tabular}
\end{table}